% Compile with: pdflatex research_document.tex
\documentclass[11pt,a4paper]{article}

\usepackage[T1]{fontenc}
\usepackage{lmodern}
\usepackage[margin=2.5cm]{geometry}
\usepackage{hyperref}
\usepackage{booktabs}
\usepackage{enumitem}
\usepackage{amsmath,amssymb}
\usepackage{natbib}
\usepackage{xcolor}

\hypersetup{
  colorlinks=true,
  linkcolor=blue!70!black,
  citecolor=green!50!black,
  urlcolor=blue!60!black,
}

\title{\textbf{Deep Dive: Generative Molecular Design, Property Prediction,\\and Formulation Optimization for AI-Driven Coatings}}
\author{Generated Research Document}
\date{March 2026}

\begin{document}
\maketitle
\tableofcontents
\newpage

%% =======================================================================
\section{Executive Summary}
%% =======================================================================

This document provides a deep technical survey of three interconnected pillars of AI-driven multifunctional coating design: \emph{generative molecular design} (how AI proposes novel coating molecules), \emph{property prediction} (how ML models forecast coating performance), and \emph{formulation optimization} (how Bayesian and active-learning methods navigate the high-dimensional formulation space).
Together, these three capabilities form the computational engine of a modern coating discovery pipeline---generative models propose candidates, property predictors screen them in silico, and optimization algorithms select the next experiments to run.
We cover nine sub-topics in depth: variational autoencoders, diffusion models, graph-based generation, corrosion prediction, UV/optical prediction, mechanical property prediction, Bayesian optimization, active learning, and multi-objective optimization.
Each section provides mathematical foundations, benchmark results, software tools, and practical considerations for implementation in coating design workflows.

%% =======================================================================
\section{VAEs and Autoencoders for Coating Molecules}
\label{sec:vaes}
%% =======================================================================

\subsection{SMILES and SELFIES Representations}

Molecular generative models require a machine-readable representation.
SMILES (Simplified Molecular Input Line Entry System) encodes molecular graphs as character strings, but arbitrary string manipulation can produce invalid molecules.
SELFIES (Self-Referencing Embedded Strings), introduced by Krenn et al.\ (2020), guarantees that \emph{every} string in its grammar decodes to a valid molecule, achieving 100\% validity by construction.
This is critical for coating design, where generated candidates must be synthetically accessible.

For polymers, neither SMILES nor SELFIES natively captures repeat units, degree of polymerisation, or branching topology.
BigSMILES extends SMILES with stochastic objects (\texttt{[\$]..[\$]}) to describe polymer ensembles, while graph-based representations encode monomers as nodes and bonds as edges.
A recent large-scale SELFIES Transformer VAE (2025) introduces a loss aligning latent-space distances with chemical similarity, yielding a highly structured, chemistry-aware latent space trained on 1.5 million compounds.

\subsection{Latent Space Navigation and Property Optimization}

The VAE encodes molecules into a continuous latent vector $\mathbf{z} \in \mathbb{R}^d$ via an encoder $q_\phi(\mathbf{z}|\mathbf{x})$ and decodes via $p_\theta(\mathbf{x}|\mathbf{z})$, with the training objective:
\begin{equation}
\mathcal{L} = -\mathbb{E}_{q_\phi}[\log p_\theta(\mathbf{x}|\mathbf{z})] + \beta \, D_\text{KL}(q_\phi(\mathbf{z}|\mathbf{x}) \| p(\mathbf{z}))
\end{equation}
where $\beta$ controls the trade-off between reconstruction fidelity and latent-space regularity.

Property optimization proceeds by training a property predictor $f(\mathbf{z}) \to y$ in latent space (e.g., predicting $T_g$, UV absorption, or corrosion resistance) and then performing gradient ascent:
$\mathbf{z}_{t+1} = \mathbf{z}_t + \eta \nabla_{\mathbf{z}} f(\mathbf{z}_t)$
to navigate toward high-performance regions.
Gómez-Bombarelli et al.\ (2018) demonstrated this on 250\,k drug-like molecules, optimising $\log P$ and QED.

\subsection{Conditional VAEs for Target Properties}

Conditional VAEs (CVAEs) augment the latent space with a target property vector $\mathbf{c}$ (e.g., $T_g = 150\,°$C, $\lambda_\text{max} = 350\,$nm):
\begin{equation}
p_\theta(\mathbf{x}|\mathbf{z}, \mathbf{c}), \quad q_\phi(\mathbf{z}|\mathbf{x}, \mathbf{c})
\end{equation}
At generation time, one specifies $\mathbf{c}$ and samples $\mathbf{z} \sim p(\mathbf{z})$, producing molecules conditioned on the desired property envelope.
Property-guided generation of complex polymer topologies using CVAEs (npj Computational Materials, 2024) demonstrated design of linear, cyclic, branch, comb, star, and dendritic polymers from a dataset of 1\,342 structures.

%% =======================================================================
\section{Diffusion Models for Molecular Generation}
\label{sec:diffusion}
%% =======================================================================

\subsection{E(3)-Equivariant Diffusion Models}

The E(3) equivariant diffusion model (EDM) by Hoogeboom et al.\ (ICML 2022) learns a denoising process over 3D molecular coordinates and atom types jointly.
The forward process adds Gaussian noise to atom positions $\mathbf{x}_t = \sqrt{\bar\alpha_t}\,\mathbf{x}_0 + \sqrt{1-\bar\alpha_t}\,\boldsymbol\epsilon$, where $\boldsymbol\epsilon \sim \mathcal{N}(0, \mathbf{I})$.
The denoising network $\epsilon_\theta(\mathbf{x}_t, t)$ is constrained to be E(3)-equivariant: for any rotation $R$, translation $\mathbf{t}$, and reflection:
\begin{equation}
\epsilon_\theta(R\mathbf{x}_t + \mathbf{t}, t) = R\,\epsilon_\theta(\mathbf{x}_t, t)
\end{equation}
This ensures generated molecules respect physical symmetries.
On QM9 (134\,k molecules, up to 9 heavy atoms), EDM achieves atom stability $> 98\%$ and molecule stability $> 82\%$.

\subsection{Score-Based SDEs and Noise Schedules}

The score-based framework (Song et al., 2021) formulates diffusion as a stochastic differential equation:
\begin{equation}
d\mathbf{x} = \mathbf{f}(\mathbf{x}, t)\,dt + g(t)\,d\mathbf{w}
\end{equation}
and trains a score network $\mathbf{s}_\theta(\mathbf{x}, t) \approx \nabla_{\mathbf{x}} \log p_t(\mathbf{x})$.
The EMDS framework (BMC Bioinformatics, 2024) decomposes diffusion into separate processes for atom coordinates, atom types, and bond types, while capturing cross-component dependencies.
This separation allows different noise schedules for continuous (coordinates) and categorical (types) features---a variance-preserving schedule for coordinates and a uniform-noise schedule for types.

Noise schedule design critically affects sample quality.
Cosine schedules, linear schedules, and learned schedules have been compared; for molecular generation, cosine schedules with signal-to-noise ratio (SNR) weighting (Min-SNR-5) consistently outperform linear schedules.

\subsection{Multiscale Diffusion for Coatings}

Multiscale graph equivariant diffusion (MD3MD, Science Advances 2025) partitions molecular conformations into multiscale graphs, assigning different interaction weights at each scale.
For coating design, a coarse-to-fine approach is particularly powerful: first generate a molecular skeleton (chromophore backbone, cross-linking topology), then refine to full atomic detail.
This mirrors the physical hierarchy of coating function: macroscopic properties (film thickness, roughness) $\to$ mesoscale (domain morphology) $\to$ molecular (monomer structure).

%% =======================================================================
\section{Graph-Based Generation and Inverse Design}
\label{sec:graphgen}
%% =======================================================================

\subsection{Junction Tree VAE}

The Junction Tree VAE (JT-VAE, Jin et al., ICML 2018) generates molecules in two phases:
\begin{enumerate}[nosep]
  \item \textbf{Tree generation}: produce a tree-structured scaffold over chemically valid substructures (rings, functional groups).
  \item \textbf{Graph assembly}: combine substructures into a complete molecular graph via message-passing.
\end{enumerate}
JT-VAE achieves 100\% chemical validity by construction, versus $\sim$43\% for character-level VAEs and $\sim$81\% for grammar VAEs.
The fragment vocabulary is extracted from the training set using BRICS decomposition or ring/chain fragmentation.
For coating polymers, the vocabulary can be customised to include UV-absorbing chromophores (benzotriazole, hydroxybenzophenone), cross-linking groups (epoxide, isocyanate), and hydrophobic tails (fluoroalkyl, siloxane).

\subsection{Fragment Assembly and Scaffold Hopping}

FRATTVAE (Communications Chemistry, 2025) treats molecules as tree structures with fragments as nodes, enabling efficient handling of large and complex compounds.
Fragment-based methods reduce the generation graph size by 5--10$\times$ compared to atom-level models, shortening training and generation times proportionally.

Scaffold hopping---replacing the core scaffold while preserving peripheral functional groups---is a key operation for coating design.
Given a known UV absorber, scaffold hopping can propose novel cores with improved photostability or altered absorption wavelength while retaining the same pendant groups for cross-linking compatibility.

\subsection{Inverse Design with Reinforcement Learning}

Inverse design pipelines specify target properties and search for molecules that satisfy them.
RL-guided generation uses a policy network (often a GNN or RNN) that proposes molecular edits, with a reward function combining property predictions:
\begin{equation}
R(\mathbf{x}) = w_1 f_\text{UV}(\mathbf{x}) + w_2 f_\text{corr}(\mathbf{x}) + w_3 f_\text{synth}(\mathbf{x})
\end{equation}
where $f_\text{UV}$ is the predicted UV absorption quality, $f_\text{corr}$ is corrosion resistance, and $f_\text{synth}$ is a synthetic accessibility score.
Ligase-conditioned JT-VAE (LC-JT-VAE) has demonstrated conditional generation of molecules with targeted properties in constrained chemical spaces.

%% =======================================================================
\section{Corrosion Resistance Prediction}
\label{sec:corrosion}
%% =======================================================================

\subsection{EIS Modeling with Machine Learning}

Electrochemical impedance spectroscopy (EIS) measures the complex impedance $Z(\omega) = Z'(\omega) + jZ''(\omega)$ of a coated substrate as a function of frequency.
The low-frequency impedance modulus $|Z|_{0.01\,\text{Hz}}$ is a widely used proxy for barrier quality.
ML models map formulation parameters (binder type, filler loading, curing schedule) to $|Z|$ values.

The ZIF-8@Ca self-healing coating study (npj Materials Degradation, 2024) used the orthogonal Latin-square method to design an initial 16-experiment matrix varying:
\begin{itemize}[nosep]
  \item Polyetheramine molecular weight (230, 400, 2000 Da)
  \item Molar ratio of polyetheramine to epoxy
  \item UPy-D400 hydrogen-bond unit content
  \item ZIF-8@Ca microfiller mass content
\end{itemize}
A random forest model was trained on these 16 points and used for active learning.
After 30 total experiments, the optimised coating achieved $|Z| > 10^9\,\Omega\cdot\text{cm}^2$---an order of magnitude above the initial best.

\subsection{Salt Spray and Accelerated Weathering Models}

Salt spray testing (ASTM B117) reports hours to first visible corrosion.
XGBoost and gradient-boosted tree models trained on datasets of 200--500 coating formulations predict hours-to-failure with $R^2 \approx 0.80$--$0.90$.
Key features include dry film thickness ($\mu$m), pigment volume concentration (\%), binder $T_g$, and inhibitor loading.

\subsection{Degradation Forecasting}

Long-term degradation follows non-linear trajectories influenced by temperature, humidity, UV exposure, and salt concentration.
Time-series models (LSTM, temporal convolutional networks) trained on accelerated-weathering EIS time traces can forecast impedance decline, though extrapolation beyond the training window remains unreliable.
Physics-informed neural networks (PINNs) that embed Fick's diffusion law for water uptake show promise for extending prediction horizons.

%% =======================================================================
\section{UV and Optical Property Prediction}
\label{sec:uvpred}
%% =======================================================================

\subsection{TD-DFT Surrogate Models}

Time-dependent density functional theory (TD-DFT) computes excited-state energies and oscillator strengths but is expensive ($\mathcal{O}(N^3)$ to $\mathcal{O}(N^5)$ scaling).
ML surrogates trained on TD-DFT data (10\,k--100\,k molecules) predict the $S_0 \to S_1$ excitation energy with mean absolute errors of 0.15--0.25\,eV, sufficient to screen candidates into UVA (315--400\,nm), UVB (280--315\,nm), or UVC ($<280$\,nm) windows.
SchNet and DimeNet++ achieve the best accuracy on the QM9 excited-state subset, while transfer learning from ground-state pre-trained models further reduces data requirements.

\subsection{GNN-Based Absorption Spectra Prediction}

Graph neural networks operating on molecular graphs predict full UV-vis absorption spectra (not just $\lambda_\text{max}$).
The input graph encodes atoms as nodes (features: element, hybridisation, formal charge) and bonds as edges (features: bond order, conjugation, ring membership).
Message-passing layers aggregate local chemical environments, and a readout layer produces the spectrum as a vector of absorbance values at discrete wavelengths.
For coating chromophores, models trained on databases of organic dyes (e.g., the DSSC dye database, $\sim$5\,000 entries) achieve spectral-shape correlation $> 0.90$.

\subsection{Photostability Modeling}

Photostability---resistance to UV-induced degradation---requires understanding of triplet-state dynamics, intersystem crossing rates, and radical formation.
These are computationally expensive to simulate (multi-reference methods, spin-orbit coupling).
Transfer-learning approaches first pre-train on abundant singlet-state data, then fine-tune on scarce triplet-state data (typically 500--2\,000 molecules).
A practical proxy is the singlet-triplet energy gap $\Delta E_{ST}$: molecules with large gaps ($> 1$\,eV) tend to be more photostable because intersystem crossing is kinetically suppressed.

%% =======================================================================
\section{Mechanical Property Prediction}
\label{sec:mechpred}
%% =======================================================================

\subsection{Glass Transition and Melting Temperature}

Graph convolutional networks (GCNs) predict $T_g$ and $T_m$ for polymers with reliable accuracy.
Tao et al.\ (ACS Polymers Au, 2021) trained GCN models on datasets of 600--1\,200 polymers, achieving $R^2 > 0.85$ for both $T_g$ and $T_m$.
Polymer representations use repeat-unit graphs with features encoding element identity, hybridisation, and degree.
The CRIPT (Community Resource for Innovation in Polymer Technology) database consortium is expanding training data availability.

\subsection{Adhesion, Hardness, and Multitask GNNs}

PolymerGNN (npj Computational Materials, 2023) is a multitask architecture that simultaneously predicts multiple properties---$T_g$, density, elastic modulus, and refractive index---from a shared graph representation.
Multitask learning improves prediction for data-scarce properties (e.g., adhesion energy, with $< 200$ labelled examples) by transferring knowledge from data-rich tasks.
Feature extraction is 1--2 orders of magnitude faster than handcrafted molecular descriptors (e.g., RDKit fingerprints), enabling rapid virtual screening.

Polymer Informatics at Scale (Chemistry of Materials, 2023) demonstrated polymer GNN models trained on 15\,000+ data points across 8 property tasks, with consistent improvements from multi-task learning versus single-task baselines.

\subsection{Elastic Modulus and MD Validation}

Elastic modulus predictions from GNN models are validated against molecular dynamics (MD) simulations.
All-atom MD with force fields like OPLS-AA or GAFF computes stress-strain curves, yielding Young's modulus directly.
For high-throughput screening, coarse-grained MD (MARTINI, SDK) reduces simulation time by 100$\times$ at the cost of some accuracy.
ML-MD hybrid workflows use GNN predictions to pre-screen candidates, then run full MD on the top 50--100 for validation.

%% =======================================================================
\section{Bayesian Optimization for Coating Formulations}
\label{sec:bayesopt}
%% =======================================================================

\subsection{Gaussian Process Surrogates}

A Gaussian process (GP) defines a distribution over functions:
\begin{equation}
f(\mathbf{x}) \sim \mathcal{GP}(m(\mathbf{x}), k(\mathbf{x}, \mathbf{x}'))
\end{equation}
where $m(\mathbf{x})$ is the mean function and $k(\mathbf{x}, \mathbf{x}')$ is the kernel.
For coating formulations with heterogeneous variables (wt\% filler, curing temperature, binder type), anisotropic kernels (e.g., ARD Mat\'ern 5/2) with per-dimension length scales outperform isotropic kernels.
Benchmarking across 14 materials science tasks showed ARD kernels and random forest surrogates achieve comparable performance, both significantly outperforming isotropic GPs.

Mixed-variable GPs handle categorical inputs (binder chemistry, solvent type) via one-hot encoding or Gower distance kernels, enabling BO over the full coating formulation space.

\subsection{Acquisition Functions}

The acquisition function balances exploration (regions of high uncertainty) and exploitation (regions of high predicted performance):
\begin{itemize}[nosep]
  \item \textbf{Expected Improvement (EI)}: $\text{EI}(\mathbf{x}) = \mathbb{E}[\max(0, f(\mathbf{x}) - f^*)]$ where $f^*$ is the current best.
  \item \textbf{Upper Confidence Bound (UCB)}: $\text{UCB}(\mathbf{x}) = \mu(\mathbf{x}) + \kappa\,\sigma(\mathbf{x})$ with exploration parameter $\kappa$.
  \item \textbf{Knowledge Gradient (KG)}: measures the expected improvement in the \emph{optimal decision} after observing one more sample; particularly effective for noisy coating measurements.
\end{itemize}
For batch experimentation (common in coating labs where 10--20 samples are prepared simultaneously), $q$-EI and $q$-UCB extensions select batches that collectively maximise information gain.

\subsection{Software Ecosystem}

\begin{table}[h]
\centering
\caption{Key Bayesian optimization software tools for coatings.}
\label{tab:software}
\small
\begin{tabular}{@{}llll@{}}
\toprule
\textbf{Package} & \textbf{Backend} & \textbf{Key Feature} & \textbf{License} \\
\midrule
BoTorch & GPyTorch / PyTorch & Batch MOBO, MC acquisition & MIT \\
BayesianOptimization & scikit-learn & Simple API, 1-D to 20-D & MIT \\
ProcessOptimizer & scikit-optimize & Experimentalist-friendly & BSD-3 \\
Ax Platform & BoTorch & Industrial-scale trials & MIT \\
pymoo & NumPy & NSGA-II/III, ref.\ directions & Apache 2.0 \\
\bottomrule
\end{tabular}
\end{table}

BoTorch (Meta) provides the most flexible framework, supporting custom GP models, multi-output surrogates, and Monte Carlo acquisition functions.
ProcessOptimizer (JCIM, 2025) specifically targets experimentalist scientists in chemistry and coatings, offering built-in plotting, noise handling, and multi-objective support with minimal configuration.

%% =======================================================================
\section{Active Learning Strategies}
\label{sec:activelearn}
%% =======================================================================

\subsection{Query Strategies for Coating Data}

Active learning iteratively selects the most informative experiment to run next:
\begin{itemize}[nosep]
  \item \textbf{Uncertainty sampling}: select the point $\mathbf{x}^* = \arg\max_{\mathbf{x}} \sigma(\mathbf{x})$ where the surrogate model is most uncertain.
  \item \textbf{Expected improvement}: same as BO's EI, combining uncertainty with proximity to the optimum.
  \item \textbf{Query-by-committee (QBC)}: train an ensemble of models and select the point with highest disagreement.
\end{itemize}
Lookman et al.\ (npj Computational Materials, 2019) compared these strategies on materials discovery tasks, finding that EI and uncertainty sampling achieve near-optimal performance across diverse design spaces, while QBC is more robust when the surrogate model is misspecified.

\subsection{Sample Efficiency Gains}

Active learning typically reduces the number of experiments by 3--5$\times$ versus random or grid-based sampling.
In the ZIF-8@Ca coating study, 30 actively selected experiments outperformed 16 Latin-square experiments plus 100 random samples.
For coating formulations with 5--15 design variables, active learning becomes essential: random sampling of a 10-dimensional space at 5 levels requires $5^{10} \approx 10$ million experiments, while active learning converges in 50--200.

Data augmentation strategies---augmenting experimental data with physics-based simulations (MD, DFT) or transfer from related coating systems---further improve sample efficiency.
Multi-fidelity active learning combines cheap low-fidelity data (e.g., DFT with a small basis set) with expensive high-fidelity data (e.g., experimental EIS) using co-kriging surrogates.

\subsection{Closed-Loop Integration}

In a closed-loop active learning system, the AI proposes an experiment, a robotic platform (e.g., Polybot) executes it, characterisation data flows back into the surrogate model, and the cycle repeats.
Convergence criteria include:
\begin{itemize}[nosep]
  \item \textbf{EI threshold}: stop when $\text{EI}_\text{max} < \epsilon$ (no significant improvement expected).
  \item \textbf{Uncertainty threshold}: stop when $\max_{\mathbf{x}} \sigma(\mathbf{x}) < \delta$ (the model is confident everywhere).
  \item \textbf{Budget exhaustion}: stop after $N$ iterations (practical constraint).
\end{itemize}

%% =======================================================================
\section{Multi-Objective Optimization}
\label{sec:moopt}
%% =======================================================================

\subsection{Pareto Front Discovery}

Real coatings optimise conflicting objectives simultaneously---e.g., maximise UV absorption while minimising $T_g$ (for flexibility) and cost.
The Pareto front $\mathcal{P}$ is the set of solutions where no objective can be improved without degrading another:
\begin{equation}
\mathcal{P} = \{ \mathbf{x} \mid \nexists \mathbf{x}' : f_i(\mathbf{x}') \geq f_i(\mathbf{x}) \;\forall i, \; f_j(\mathbf{x}') > f_j(\mathbf{x}) \;\exists j \}
\end{equation}

Multi-objective evolutionary algorithms (NSGA-II, NSGA-III) use non-dominated sorting and crowding distance to maintain a diverse Pareto approximation.
For coating formulations, 2--4 objectives are typical; beyond 4 objectives, many-objective algorithms (RVEA, MOEA/D) with reference-direction decomposition are preferred.

\subsection{Constraint Handling for Regulatory Compliance}

Coating formulations face hard constraints:
\begin{itemize}[nosep]
  \item \textbf{VOC limits}: EU Directive 2004/42/EC limits VOC content to 200--500\,g/L depending on category.
  \item \textbf{REACH restrictions}: SVHC substances (e.g., certain isocyanates, chromium compounds) face authorisation requirements.
  \item \textbf{PFAS phase-out}: US EPA and EU ECHA are restricting per- and polyfluoroalkyl substances in coatings.
\end{itemize}
Constrained MOBO incorporates these as inequality constraints in the acquisition function, either via feasibility-weighted EI or through penalty methods.
The constraint-aware approach prunes infeasible regions before any experiment is run, saving resources.

\subsection{MOBO Frameworks and Implementation}

BoTorch provides \texttt{qNoisyExpectedHypervolumeImprovement} (qNEHVI) as its primary MOBO acquisition function, measuring the expected improvement in the hypervolume enclosed by the Pareto front.
pymoo implements NSGA-II, NSGA-III, and MOEA/D with Python bindings, supporting custom evaluation functions that can wrap ML property predictors.

A practical MOBO workflow for coating optimisation:
\begin{enumerate}[nosep]
  \item Define 2--4 objectives (corrosion resistance, UV absorption, flexibility, cost).
  \item Define constraints (VOC $< 250$\,g/L, no SVHC substances).
  \item Fit a multi-output GP on initial data (20--50 experiments).
  \item Iterate: propose batch via qNEHVI, execute experiments, update GP.
  \item After convergence, present the Pareto front to domain experts for trade-off selection.
\end{enumerate}

%% =======================================================================
\section{Related Work and State of the Art}
\label{sec:related}
%% =======================================================================

\begin{table}[h]
\centering
\caption{Comparison of generative, predictive, and optimization approaches for coatings (2022--2026).}
\label{tab:comparison}
\small
\begin{tabular}{@{}p{3.2cm}p{2.5cm}p{3cm}p{4cm}@{}}
\toprule
\textbf{Method} & \textbf{Category} & \textbf{Application} & \textbf{Key Metric} \\
\midrule
SELFIES Transformer VAE & Generative & Molecule design & 100\% validity, 1.5M training set \\
E(3)-EDM & Generative & 3D molecular gen. & 98\% atom stability on QM9 \\
JT-VAE & Generative & Fragment-based gen. & 100\% validity, 5--10$\times$ smaller graphs \\
PolymerGNN & Predictive & Multi-property & $R^2 > 0.85$ for $T_g$, $T_m$, $\rho$ \\
RF + Active Learning & Optimisation & Self-healing epoxy & $|Z| > 10^9$ in 30 experiments \\
BoTorch qNEHVI & Optimisation & Multi-objective & Hypervolume-optimal Pareto fronts \\
ProcessOptimizer & Optimisation & PU lacquer & Near-optimal in 50 experiments \\
\bottomrule
\end{tabular}
\end{table}

\paragraph{Gaps.}
(1)~No end-to-end pipeline yet combines molecular generation $\to$ property prediction $\to$ BO-driven formulation $\to$ robotic validation in a single closed loop for coatings.
(2)~Generative models are trained predominantly on small-molecule datasets (QM9, ZINC); polymer-specific generative models remain scarce.
(3)~Multi-objective BO for coatings rarely handles more than 3 objectives simultaneously, despite real formulations having 5+ performance criteria.
(4)~Explainability: SHAP and attention-based interpretability methods are underexplored in coating ML models.

%% =======================================================================
\section{Cross-Cutting Connections}
\label{sec:crosscut}
%% =======================================================================

\paragraph{VAEs $\leftrightarrow$ Diffusion Models.}
Both encode molecules in continuous spaces but with different training objectives (reconstruction + KL vs.\ denoising score matching).
Hybrid architectures---latent diffusion models that run diffusion in the VAE's latent space rather than in data space---combine the structural guarantees of VAEs with the sample quality of diffusion, and are beginning to be explored for molecular generation.

\paragraph{Graph Generation $\leftrightarrow$ Corrosion Prediction.}
Every molecule proposed by a graph-based generator must be screened for corrosion resistance before synthesis.
GNN property predictors can be called as fast oracles ($\sim$1\,ms per molecule) within the generation loop, enabling on-the-fly rejection of candidates with predicted $|Z| < 10^7\,\Omega\cdot\text{cm}^2$.

\paragraph{Bayesian Optimization $\leftrightarrow$ Corrosion / EIS.}
BO selects the next EIS experiment to maximise information gain about the impedance landscape.
The GP posterior uncertainty $\sigma(\mathbf{x})$ identifies formulation regions where the coating's corrosion resistance is most uncertain, directing experiments where they matter most.

\paragraph{UV Prediction $\leftrightarrow$ Multi-Objective Optimization.}
UV absorption is one objective in a multi-objective formulation problem (alongside corrosion resistance, flexibility, and cost).
The UV property predictor serves as one objective function in the MOBO loop, evaluated at each candidate point.

\paragraph{Active Learning $\leftrightarrow$ Mechanical Prediction.}
Mechanical properties (adhesion, hardness) have the scarcest experimental data.
Active learning prioritises experiments that reduce uncertainty in the mechanical-property surrogate model, improving the weakest link in the multi-property prediction chain.

%% =======================================================================
\section{Key References and Resources}
\label{sec:references}
%% =======================================================================

\subsection{Papers}
\begin{itemize}[nosep]
  \item Hoogeboom et al., ``Equivariant Diffusion for Molecule Generation in 3D,'' \textit{ICML 2022}. \url{https://arxiv.org/abs/2203.17003}
  \item Jin et al., ``Junction Tree VAE for Molecular Graph Generation,'' \textit{ICML 2018}. \url{https://arxiv.org/abs/1802.04364}
  \item Park et al., ``Polymer GNN for Multitask Property Learning,'' \textit{npj Comp.\ Mat.} (2023). \url{https://www.nature.com/articles/s41524-023-01034-3}
  \item Lookman et al., ``Active Learning in Materials Science,'' \textit{npj Comp.\ Mat.} (2019). \url{https://www.nature.com/articles/s41524-019-0153-8}
  \item Chen et al., ``ML-Assisted Self-Healing Epoxy Coating,'' \textit{npj Mat.\ Deg.} (2024). \url{https://www.nature.com/articles/s41529-024-00427-z}
  \item ProcessOptimizer, \textit{JCIM} (2025). \url{https://pubs.acs.org/doi/10.1021/acs.jcim.4c02240}
  \item Xu et al., ``MD3MD: Multiscale Graph Equivariant Diffusion,'' \textit{Science Advances} (2025). \url{https://www.science.org/doi/10.1126/sciadv.adv0778}
  \item Krenn et al., ``SELFIES: A Robust Representation of Semantically Constrained Graphs,'' \textit{Mach.\ Learn.: Sci.\ Technol.} (2020).
  \item FRATTVAE, ``Leveraging Tree-Transformer VAE with Fragment Tokenization,'' \textit{Comm.\ Chem.} (2025). \url{https://www.nature.com/articles/s42004-025-01640-w}
\end{itemize}

\subsection{Code Repositories}
\begin{itemize}[nosep]
  \item E(3) Diffusion: \url{https://github.com/ehoogeboom/e3_diffusion_for_molecules}
  \item JT-VAE: \url{https://github.com/wengong-jin/icml18-jtnn}
  \item BoTorch: \url{https://botorch.org}
  \item pymoo: \url{https://pymoo.org}
  \item ProcessOptimizer: \url{https://pypi.org/project/ProcessOptimizer/}
  \item BayesianOptimization: \url{https://github.com/bayesian-optimization/BayesianOptimization}
  \item Molecular design papers list: \url{https://github.com/AspirinCode/papers-for-molecular-design-using-DL}
\end{itemize}

\subsection{Datasets}
\begin{itemize}[nosep]
  \item QM9: 134\,k molecules, 19 DFT properties. \url{https://moleculenet.org}
  \item ZINC 250k: drug-like molecules for generative model benchmarking.
  \item GEOM-Drugs: 450\,k conformers of drug-like molecules.
  \item CRIPT polymer database: \url{https://criptapp.org}
  \item Materials Project: 150\,k+ inorganic structures. \url{https://materialsproject.org}
\end{itemize}

%% =======================================================================
\section{Implementation Roadmap}
\label{sec:roadmap}
%% =======================================================================

\subsection{Phase 1: Data Infrastructure (Months 1--3)}
\begin{enumerate}[nosep]
  \item Curate a coating-specific molecular dataset: extract UV-absorber, corrosion-inhibitor, and polymer binder structures from literature and patents (target: 5\,000+ structures with measured properties).
  \item Standardise data schema using BigSMILES for polymers and SELFIES for small molecules.
  \item Set up a PostgreSQL + RDKit database with property annotations ($T_g$, $|Z|$, $\lambda_\text{max}$, adhesion).
\end{enumerate}

\subsection{Phase 2: Property Prediction Models (Months 2--5)}
\begin{enumerate}[nosep]
  \item Train GNN models (GCN or MPNN) on curated data for $T_g$, corrosion resistance ($|Z|$), and UV absorption ($\lambda_\text{max}$).
  \item Benchmark against random forest and XGBoost baselines.
  \item Fine-tune MACE-MP foundation model for polymer-specific force-field predictions.
  \item Validate top candidates with MD simulations (OPLS-AA) and TD-DFT (B3LYP/6-31G*).
\end{enumerate}

\subsection{Phase 3: Generative Model Training (Months 3--7)}
\begin{enumerate}[nosep]
  \item Train a JT-VAE with a coating-specific fragment vocabulary (chromophores, cross-linkers, hydrophobic tails).
  \item Train an E(3)-equivariant diffusion model on GEOM-Drugs + coating augmentation data.
  \item Implement inverse design loop: CVAE conditioned on target property vector $\to$ GNN screening $\to$ top-$k$ selection.
  \item Benchmark novelty, validity, and predicted-property distributions.
\end{enumerate}

\subsection{Phase 4: Optimization Pipeline (Months 5--9)}
\begin{enumerate}[nosep]
  \item Implement MOBO loop using BoTorch with qNEHVI acquisition function.
  \item Define 3 objectives (UV absorption, corrosion resistance, flexibility) + VOC constraint.
  \item Run initial 30 experiments (formulation + spin-coating + characterisation).
  \item Deploy active learning (uncertainty sampling) to select subsequent batches of 10.
  \item Iterate until Pareto front stabilises ($<5\%$ hypervolume change over 3 iterations).
\end{enumerate}

\subsection{Phase 5: Integration and Validation (Months 8--12)}
\begin{enumerate}[nosep]
  \item Connect generative model $\to$ property predictor $\to$ BO formulation optimizer in a single pipeline.
  \item Validate top 10 Pareto-optimal formulations with full characterisation suite (EIS, UV-vis, nanoindentation, salt spray, QUV accelerated weathering).
  \item Compare AI-designed coatings against commercial benchmarks (Tinuvin, Irganox).
  \item Publish results and open-source the pipeline.
\end{enumerate}

\subsection{Key Dependencies and Milestones}

\begin{table}[h]
\centering
\caption{Implementation milestones and dependencies.}
\label{tab:milestones}
\small
\begin{tabular}{@{}lll@{}}
\toprule
\textbf{Milestone} & \textbf{Target} & \textbf{Depends On} \\
\midrule
Curated dataset ($>$5k structures) & Month 3 & Literature mining \\
GNN property models ($R^2 > 0.80$) & Month 5 & Dataset \\
JT-VAE + EDM trained & Month 7 & Dataset \\
MOBO Pareto front (3 objectives) & Month 9 & Property models \\
Validated formulations ($n=10$) & Month 12 & Full pipeline \\
\bottomrule
\end{tabular}
\end{table}

%% =======================================================================
\section{Conclusion}
\label{sec:conclusion}
%% =======================================================================

The three pillars examined here---generative molecular design, property prediction, and formulation optimization---form the computational core of next-generation coating discovery.
Generative models (VAEs, diffusion, graph-based) can propose millions of novel coating molecules; GNN-based property predictors screen them in milliseconds; and Bayesian/active-learning optimizers navigate the formulation landscape in tens of experiments rather than thousands.

Key open opportunities at the intersection of these pillars include:
\begin{enumerate}[nosep]
  \item \textbf{End-to-end differentiable pipelines}: backpropagating property gradients through the generative model to directly optimise for target coating performance.
  \item \textbf{Polymer-specific generative models}: extending JT-VAE and diffusion models to handle repeat units, degree of polymerisation, and cross-linking topology natively.
  \item \textbf{Multi-fidelity active learning}: combining DFT, MD, and experimental data at different fidelity levels in a unified GP surrogate to maximise information per dollar spent.
  \item \textbf{Explainable property prediction}: using attention maps and SHAP values to identify which molecular substructures drive corrosion resistance, UV absorption, or adhesion---enabling human-AI collaborative design.
  \item \textbf{Regulatory-aware MOBO}: dynamically incorporating evolving PFAS, VOC, and REACH constraints as the regulatory landscape shifts, ensuring generated formulations remain compliant.
\end{enumerate}

The convergence of these methods, combined with autonomous experimentation platforms, promises to compress the coating discovery cycle from years to weeks.

\end{document}
